% Farbpalette der Layer/Bahnen -- einmal in die Präambel, dann in Text, Diagrammen und Plots nutzbar.
% Kräftig = Bahnkopf und Plot-Kurven | Rand = Bahnrahmen | Pastell = Blockfüllung | Block = Blockrand
% Benötigt: \usepackage{xcolor}
\definecolor{FrontendKraeftig}{HTML}{4A90D9}  \definecolor{FrontendRand}{HTML}{2F6FB0}  \definecolor{FrontendPastell}{HTML}{DAE8FC}  \definecolor{FrontendBlock}{HTML}{6C8EBF}
\definecolor{BackendKraeftig}{HTML}{9A9A9A}   \definecolor{BackendRand}{HTML}{5F5F5F}   \definecolor{BackendPastell}{HTML}{F5F5F5}   \definecolor{BackendBlock}{HTML}{666666}
\definecolor{LayerEinsKraeftig}{HTML}{5CB85C} \definecolor{LayerEinsRand}{HTML}{3E8E41} \definecolor{LayerEinsPastell}{HTML}{D5E8D4} \definecolor{LayerEinsBlock}{HTML}{82B366}  % Pi-Decke
\definecolor{LayerZweiKraeftig}{HTML}{F2B632} \definecolor{LayerZweiRand}{HTML}{C48A00} \definecolor{LayerZweiPastell}{HTML}{FFF2CC} \definecolor{LayerZweiBlock}{HTML}{D6B656}  % Pi-Hand
\definecolor{RoboterKraeftig}{HTML}{F28C28}   \definecolor{RoboterRand}{HTML}{C25E00}   \definecolor{RoboterPastell}{HTML}{FFE6CC}   \definecolor{RoboterBlock}{HTML}{D79B00}
